
% !TEX TS-program = xelatex
\documentclass[11pt]{article}

% -----------------------------
% Packages
% -----------------------------
\usepackage[a4paper,margin=1in]{geometry}
\usepackage{fontspec}
\usepackage{microtype}
\usepackage{hyperref}
\usepackage{bookmark}
\usepackage{longtable}
\usepackage{array}
\usepackage{booktabs}
\usepackage{xcolor}
\usepackage{enumitem}
\usepackage{titlesec}
\usepackage{fancyhdr}
\usepackage{graphicx}
\usepackage{listings}
\usepackage{amsmath,amssymb}
\usepackage{verbatim}
\usepackage{caption}

\hypersetup{
  colorlinks=true,
  linkcolor=blue,
  urlcolor=blue,
  citecolor=blue
}

\setmainfont{TeX Gyre Pagella}

% -----------------------------
% Styling
% -----------------------------
\titleformat{\section}{\Large\bfseries}{\thesection}{0.75em}{}
\titleformat{\subsection}{\large\bfseries}{\thesubsection}{0.75em}{}
\titleformat{\subsubsection}{\normalsize\bfseries}{\thesubsubsection}{0.75em}{}

\setlist[itemize]{topsep=4pt,itemsep=3pt,parsep=0pt,partopsep=0pt}

\pagestyle{fancy}
\fancyhf{}
\lhead{Sovereign Desktop Superintelligence}
\rhead{Local-First Formal-Agentic Engineering}
\cfoot{\thepage}

% -----------------------------
% Document metadata
% -----------------------------
\title{\textbf{Sovereign Desktop Superintelligence}\\
\large Local-First Formal-Agentic Engineering Pipeline\\
\vspace{0.25em}
\small (Includes: Recommended Full Local Stack --- March 2026)}
\author{Dylan Kawalec\\\small chaotic beauty LLC}
\date{March 5, 2026\\\small Version 1.0 \quad\textbar\quad License: CC BY 4.0 (recommended)}

\begin{document}
\maketitle

\begin{abstract}
This document presents a complete, closed-loop, local-first engineering factory that generates
reviewable artifacts (specs, invariants, proofs/checks, code, tests, builds) instead of opaque outputs.
It begins from human-authored invariants, formalizes behavior with TLA+/PlusCal and verification-aware
languages (e.g., Dafny), compiles verified requirements into a living PRD, and drives agentic code synthesis
and repair using locally hosted open-weight models. Observability and runtime isolation (eBPF, OpenTelemetry,
WebAssembly/Tauri) are used to create traceable telemetry and to distill recurring patterns into Executable Service
Components (ESCs). Any external “sanity oracle” (e.g., Grok via API) is explicitly optional and out-of-trust-boundary.
\end{abstract}

\tableofcontents
\newpage

% ============================================================
\section{Executive Summary}

The goal is to build a desktop ``AI engineering OS'' that:
\begin{itemize}
  \item Runs primarily on local hardware with open-source tools and open-weight models.
  \item Is composable and inspectable end-to-end.
  \item Produces \emph{reviewable} artifacts (specs, proofs, code, tests, builds) rather than unverifiable outputs.
  \item Enforces traceability from every generated artifact back to human-authored invariants.
\end{itemize}

The system is designed as a deterministic loop with explicit boundaries:
\begin{itemize}
  \item \textbf{In-trust:} local inference, local orchestration, formal checking, local testing, local artifact storage.
  \item \textbf{Out-of-trust (optional):} any external oracle/model API calls used only for critique/triangulation.
\end{itemize}

% ============================================================
\section{Core Layers of the Pipeline}

\subsection{Human Invariant Layer}

The Human Invariant Layer is the sole human-only step where the engineer authors the immutable foundation
of correctness and intent, using the \textbf{APCL methodology}:

\begin{itemize}
  \item \textbf{Axioms:} non-negotiable truths and constraints.
  \item \textbf{Postulates:} assumed properties and operating conditions.
  \item \textbf{Corollaries:} derived constraints that must hold if axioms/postulates hold.
  \item \textbf{Lemmas:} intermediate claims required to prove corollaries or discharge proof obligations.
\end{itemize}

The invariants also include:
\begin{itemize}
  \item Safety and liveness properties (temporal logic).
  \item Resource budgets (token limits, vCPU cores, RAM, disk, latency caps).
  \item Loop intent (purpose, end-state, success metrics, termination and escalation rules).
\end{itemize}

\subsection{OODA Execution Layer (Local Open-Source Assistants)}

The OODA layer is the persistent autonomous ``engine room'' of the workflow.
OODA stands for \textbf{Observe -- Orient -- Decide -- Act}.

\subsubsection{Inputs}
Each iteration ingests:
\begin{itemize}
  \item Research artifacts (papers, notes, PDFs, Markdown).
  \item Codebase state (git diffs, issues, history).
  \item Exhaustive test results (unit, integration, fuzz/property tests).
  \item Deep telemetry (eBPF probes, OpenTelemetry traces/metrics/logs).
\end{itemize}

\subsubsection{Model/Tooling Examples}
Examples of local assistants include:
\begin{itemize}
  \item Terminal-native multi-agent coding assistants (e.g., CodeBuff-like role).
  \item Local forks of large coding models (e.g., Kimi K2.5).
  \item DeepSeek-Coder-style families (e.g., DeepSeek-V3.2) for long-horizon engineering.
\end{itemize}

\subsubsection{OODA Stages}
\begin{enumerate}[leftmargin=2em]
  \item \textbf{Observe:} watch folders, git hooks, test logs, telemetry streams.
  \item \textbf{Orient:} cross-reference observations against APCL invariants and long-term memory.
  \item \textbf{Decide:} produce a concrete action plan (regenerate validator, strengthen tests, distill telemetry).
  \item \textbf{Act:} execute changes, compile/regenerate micro-components, update memory, log traceability.
\end{enumerate}

\subsubsection{Hypervelocity Components and ESCs}
\textbf{Hypervelocity components} are ultra-lightweight single-responsibility skill binaries (often $<5$ MB)
regenerated on-demand to remain tailored to the current task.

\textbf{Executable Service Components (ESCs)} are self-learning service binaries baked into the local environment,
distilling recurring runtime patterns into optimized components.

\subsection{Formal Specification \& Verification Layer}

This layer converts the refined understanding from OODA into a mathematically checkable blueprint.

\subsubsection{Key Concepts}
\begin{itemize}
  \item \textbf{TLA+}: Temporal Logic of Actions specification language for concurrent/distributed systems.
  \item \textbf{PlusCal}: algorithmic notation translated into TLA+.
  \item \textbf{TLC}: explicit-state model checker.
  \item \textbf{Apalache}: symbolic model checker for larger state spaces.
  \item \textbf{Dafny}: verification-aware language/compiler for executable specifications and proofs.
  \item \textbf{SMT solvers}: Z3/cvc5 to discharge proof obligations.
\end{itemize}

\subsubsection{Formal Process}
\begin{enumerate}[leftmargin=2em]
  \item \textbf{One-shot formal emission:} a model emits a self-contained PlusCal specification from invariants + OODA output.
  \item \textbf{Translation \& model checking:} PlusCal $\rightarrow$ TLA+; TLC/Apalache verify invariants and temporal properties.
  \item \textbf{Green-path collapse to living PRD:} when checks pass, produce a canonical \LaTeX{} PRD with traceability.
  \item \textbf{Hierarchical prompt tree (A--B--C units):} generate a strict prompt hierarchy derived from the proven spec.
\end{enumerate}

\subsection{Sanity \& PR Layer (Optional External Gate)}

This layer is a final quality gate before synthesis. It may include:
\begin{itemize}
  \item An optional external oracle critique (e.g., Grok via API) using redacted derived summaries only.
  \item Automated PR drafting (e.g., IDE composer).
  \item Traceability audits mapping diffs back to formal spec and invariants.
\end{itemize}

\textbf{Sovereignty boundary:} any external call is optional and explicitly non-sovereign.

\subsection{One-Shot Synthesis \& Commit Layer}

\begin{enumerate}[leftmargin=2em]
  \item \textbf{One-shot code synthesis:} a local model reads the ABC prompt tree and generates complete code.
  \item \textbf{Full verification execution:} run TDD suite; re-run TLC; optional lightweight critique.
  \item \textbf{Signed auto-commit \& push:} if all checks are green, create a signed commit with traceability artifacts.
\end{enumerate}

% ============================================================
\section{Recommended Full Local Stack (March 2026)}
\label{sec:stack}

\subsection{Design Goals and Non-Goals}

\subsubsection{Goals}
\begin{itemize}
  \item \textbf{Local-first:} run primarily on owned hardware via open models and local servers.
  \item \textbf{Auditable:} produce inspectable artifacts (specs, checks, proofs, traces, diffs).
  \item \textbf{Composable:} swap inference engines, models, and tools via OpenAI-compatible endpoints where possible.
  \item \textbf{Safety by construction:} isolate execution (containers/WASM), monitor behavior (eBPF/telemetry), require human sign-off.
\end{itemize}

\subsubsection{Non-Goals}
\begin{itemize}
  \item ``Perfect verification'': formal methods prove specific properties of specified models, not total correctness.
  \item ``Free inference'': local inference eliminates per-token billing but has non-zero marginal cost (power, time, amortization).
\end{itemize}

\subsection{Inference Layer (Local, Swappable)}

\begin{itemize}
  \item \textbf{Ollama} --- pull-and-run ergonomics, model management, consistent local API endpoint.
  \item \textbf{vLLM} --- high-throughput serving, batching, multi-GPU, OpenAI-compatible server.
  \item \textbf{llama.cpp} --- maximal portability, aggressive GGUF quantization, OpenAI-compatible HTTP server.
\end{itemize}

\paragraph{Hardware realism.}
Very large models (hundreds of billions of parameters) typically require multi-GPU workstations
or heavy CPU offload for usable throughput; treat 480B/744B/1T-class models as workstation/mini-rack targets.

\subsection{Model Set by Role}

\subsubsection{Specification + deep code/repo reasoning (primary, heavy)}
\begin{itemize}
  \item \textbf{Qwen3-Coder-480B-A35B-Instruct} (Apache-2.0): repo-scale reasoning and spec-to-code workflows.
  \item \textbf{Kimi K2.5} (Modified MIT): strong generalist; can be used for multimodal workflows where relevant.
\end{itemize}

\subsubsection{Agentic execution (tool use, long-horizon tasks)}
\begin{itemize}
  \item \textbf{GLM-5} (MIT): agentic-focused reasoning/coding.
  \item \textbf{DeepSeek-V3.2} (MIT): long-context + agent performance.
\end{itemize}

\subsubsection{Small ``skill models'' (cheap, fast, local-everywhere)}
Prefer distilled 7B--14B variants for linting, style transforms, unit test generation, small refactors,
classification, routing, and retrieval summarization.

\subsection{Orchestration and Multi-Agent Frameworks}

\begin{itemize}
  \item \textbf{LangGraph}: primary deterministic orchestration substrate (agents as graphs).
  \item Optional role-team templates: \textbf{MetaGPT} and \textbf{ChatDev 2.0}.
\end{itemize}

\subsection{Coding Agents and IDE Surfaces}

\begin{itemize}
  \item \textbf{OpenHands} (formerly OpenDevin): autonomous SWE workflows; note that enterprise features may not be fully sovereign.
  \item \textbf{Aider}: terminal pair-programmer supporting local Ollama and OpenAI-compatible endpoints.
  \item \textbf{Continue}: IDE assistant, usable with local backends such as Ollama.
\end{itemize}

\subsection{Formal Methods Toolchain}

\begin{itemize}
  \item \textbf{TLA+ Toolbox} and \textbf{TLC}.
  \item \textbf{Apalache}: symbolic model checking.
  \item \textbf{Dafny}: verification-aware language/compiler.
  \item \textbf{SMT solvers}: Z3 and cvc5.
\end{itemize}

\subsection{Memory and Knowledge Substrate}

\begin{itemize}
  \item \textbf{GraphRAG}: graph-based RAG pipeline with local quickstart.
  \item \textbf{Letta} (formerly MemGPT): stateful agent memory with explicit memory management.
\end{itemize}

\subsection{Runtime Isolation and Observability}

\begin{itemize}
  \item \textbf{WebAssembly security model}: sandboxed modules.
  \item \textbf{eBPF}: low-overhead kernel-level telemetry.
  \item \textbf{Tauri}: local native desktop embedding (Rust backend).
  \item \textbf{OpenTelemetry}: vendor-neutral observability standard.
\end{itemize}

\subsection{Optional External ``Sanity Oracle''}

If you choose to call an external oracle model, treat it as outside the sovereignty boundary and use it only for critique/triangulation.

% ============================================================
\section{Canonical Component Table}

\begin{longtable}{@{}p{0.11\textwidth}p{0.16\textwidth}p{0.26\textwidth}p{0.18\textwidth}p{0.12\textwidth}p{0.17\textwidth}@{}}
\caption{Canonical Component Table (links, roles, licenses, suggested starts).}\label{tab:components}\\
\toprule
\textbf{Layer} & \textbf{Component} & \textbf{Canonical link} & \textbf{Primary role} & \textbf{License} & \textbf{Suggested install/start}\\
\midrule
\endfirsthead
\toprule
\textbf{Layer} & \textbf{Component} & \textbf{Canonical link} & \textbf{Primary role} & \textbf{License} & \textbf{Suggested install/start}\\
\midrule
\endhead

Inference & Ollama & \url{https://github.com/ollama/ollama} & Local model runner + local API endpoint & MIT & \texttt{ollama run <model>}\\
Inference & vLLM & \url{https://github.com/vllm-project/vllm} & High-throughput serving + OpenAI-compat API & Apache-2.0 & \texttt{pip install vllm; vllm serve <model>}\\
Inference & llama.cpp & \url{https://github.com/ggml-org/llama.cpp} & Low-level inference + GGUF + OpenAI-compat server & MIT & Build from source; run server\\
Orchestration & LangGraph & \url{https://github.com/langchain-ai/langgraph} & Agent orchestration as graphs & MIT & \texttt{pip install -U langgraph}\\
Multi-agent & MetaGPT & \url{https://github.com/FoundationAgents/MetaGPT} & Multi-role agent workflows & MIT & \texttt{pip install --upgrade metagpt}\\
Multi-agent & ChatDev & \url{https://github.com/OpenBMB/ChatDev} & Multi-role ``software org'' template & Apache-2.0 & Follow repo instructions\\
Coding agent & OpenHands & \url{https://github.com/OpenHands/OpenHands} & Autonomous SWE agent & MIT core & Follow docs\\
Coding tool & Aider & \url{https://github.com/Aider-AI/aider} & Terminal pair-programmer & Apache-2.0 & \texttt{pip install aider-install}\\
IDE assistant & Continue & \url{https://github.com/continuedev/continue} & IDE agent + checks & Apache-2.0 & Install extension; configure\\
Formal & TLA+ & \url{https://github.com/tlaplus/tlaplus} & Spec + model checking & MIT & Install Toolbox; run TLC\\
Formal & Apalache & \url{https://github.com/apalache-mc/apalache} & Symbolic model checking & Apache-2.0 & Install Java; run \texttt{apalache-mc}\\
Formal & Dafny & \url{https://github.com/dafny-lang/dafny} & Verified programming + compilation & see repo & \texttt{dotnet tool install --global dafny}\\
SMT & Z3 & \url{https://github.com/Z3Prover/z3} & SMT solver & MIT & binaries / bindings\\
SMT & cvc5 & \url{https://github.com/cvc5/cvc5} & SMT solver & see docs & build / bindings\\
Memory & GraphRAG & \url{https://github.com/microsoft/graphrag} & Graph-based RAG pipeline & MIT & \texttt{pip install graphrag}\\
Memory & Letta & \url{https://github.com/letta-ai/letta} & Stateful memory for agents & Apache-2.0 & Run server (Docker/docs)\\
Runtime & WebAssembly & \url{https://webassembly.org/docs/security/} & Sandbox execution model & N/A & Choose runtime\\
Runtime/Obs & eBPF & \url{https://ebpf.io/what-is-ebpf/} & Kernel telemetry + safe extension & N/A & bpftool/bpftrace\\
Runtime & Tauri & \url{https://github.com/tauri-apps/tauri} & Native desktop embedding & Apache-2.0/MIT & \texttt{npm create tauri-app@latest}\\
Observability & OpenTelemetry & \url{https://opentelemetry.io/docs/} & Traces/metrics/logs standard & varies & Instrument + export OTLP\\
\bottomrule
\end{longtable}

% ============================================================
\section{Canonical Sovereign Loop (Mermaid Source)}

The following diagram is provided as Mermaid source for rendering in Mermaid-compatible tools:

\begin{lstlisting}[basicstyle=\ttfamily\small,frame=single]
flowchart TD
  A[Human intent + invariants] --> B[Spec & constraints]
  B --> C[Formalize: TLA+/Dafny properties]
  C --> D[Model check / SMT solve: TLC, Apalache, Z3/cvc5]
  D --> E[Implementation plan (agents)]
  E --> F[Agentic coding: OpenHands / Aider / Continue]
  F --> G[Build + tests + CI checks]
  G --> H[Package as ESC (local binary/app)]
  H --> I[Runtime isolation: WASM / containers / Tauri embed]
  I --> J[Observability: eBPF + OpenTelemetry traces]
  J --> K[Memory update: GraphRAG + Letta]
  K --> L{Optional external critique?}
  L -->|Yes (non-sovereign)| M[xAI Grok via API]
  L -->|No| N[Local-only]
  M --> A
  N --> A
\end{lstlisting}

% ============================================================
\section{Changelog and Corrections (Curated)}

\begin{itemize}
  \item Treat external oracles as optional and outside the local trust boundary.
  \item Prefer maintained project names (e.g., OpenHands vs.\ historical OpenDevin naming).
  \item Require explicit hardware realism notes for extremely large parameter-count models.
  \item Ensure every long-term memory item is source-linked, timestamped, and deletable.
\end{itemize}

% ============================================================
\section{Appendix: Compilation Instructions}

\subsection{Minimal build}

This source is designed to compile with \texttt{xelatex}:

\begin{lstlisting}[basicstyle=\ttfamily\small,frame=single]
xelatex sovereign-desktop-superintelligence.tex
xelatex sovereign-desktop-superintelligence.tex
\end{lstlisting}

Running twice ensures the table of contents and hyperlinks are resolved.

\end{document}
